\section*{Chapter \ref{ch:mx:dark-pools-and-blocks} --- Dark Pools and Blocks}

\begin{solution}{exo:mx:dark-pools-and-blocks:1}
The midpoint is 20.01: the buyer saves $20.02-20.01$ and the seller $20.01-20.00$, a cent a share each, USD 100 on 10\,000 shares.
\end{solution}

\begin{solution}{exo:mx:dark-pools-and-blocks:2}
Nothing with the 300-share sell (below the minimum); 1\,500 shares with the second. 3\,500 remain resting.
\end{solution}

\begin{solution}{exo:mx:dark-pools-and-blocks:3}
$28/40=70\%$. Each lapsed invitation told the participant that a large contra existed without it trading: a low rate marks a participant who uses
conditionals to fish for information, and the venue restricts it to protect the others.
\end{solution}

\begin{solution}{exo:mx:dark-pools-and-blocks:4}
The trade raises $P^\ast$ by $0.0005\times500=0.25$ ticks, 0.25 basis points at 100.00, on the 15\,000 shares still to buy: $0.25\times15\,000/20\,000=0.1875$
basis points of the order.
\end{solution}

\begin{solution}{exo:mx:dark-pools-and-blocks:5}
\omterm{def:mx:why-there-is-a-spread:informed}{Informed traders} trade when $P^\ast$ is away from the mid and all on the same side; the dark liquidity they need, on the other side, was posted by
patient uninformed traders at random, and the informed exhaust it. Uninformed takers arrive on both sides and find it more often.
\end{solution}

\begin{solution}{exo:mx:dark-pools-and-blocks:6}
Child $i$ of $q=Q/n$ is followed by $Q-iq$ shares that pay its impact $\lambda q$: $\sum_{i=1}^{n}\lambda q(Q-iq)/Q=\lambda q\,(n-(n+1)/2)=\lambda Q/2\,(1-1/n)$.
With $\lambda=0.0005$, $Q=20\,000$, $n=100$: 4.95 ticks, 4.95 basis points at 100.00.
\end{solution}

\begin{solution}{exo:mx:dark-pools-and-blocks:7}
A minimum of 200 still blocks the pinger's 100-share probes: no leakage, the whole order fills in the dark, in about 1\,020 seconds on average, at a
shortfall of 7.0 basis points (standard error 0.9) against 12.5 for the unprotected peg and 9.5 for a minimum of 1\,000; the drift is 13.6. A minimum just
above the probe size keeps the protection and loses much less of the natural flow.
\end{solution}

\begin{solution}{exo:mx:dark-pools-and-blocks:8}
Midpoint fills save the spread, not the rest: the price moved while the order was worked, because the order's fills took flow out of the lit market and
because others learnt of it. Measure the shortfall against the arrival price, the drift before completion and the counterparties' mark-outs; in the
simulation the unprotected dark order cost 12.5 basis points, all at the midpoint.
\end{solution}

\begin{solution}{pb:mx:dark-pools-and-blocks:1}
\textbf{1.} Half the spread for each side; the lit market's best bid and offer.
\textbf{2.} A current copy of the lit quotes: control N, sent after each change of the lit top.
\textbf{3.} The smallest fill a dark order accepts; it keeps small probing orders away and it reduces the fills.
\textbf{4.} A non-binding resting indication that must be firmed up when invited; the share of invitations a participant firms up.
\textbf{5.} Both sides learn that the contra exists, and one of them paid nothing to learn it.
\textbf{6.} \omterm{def:mx:why-there-is-a-spread:informed}{Informed traders} crowd on one side and face higher execution risk in the dark, so they prefer the exchange; uninformed ones prefer the dark.
\textbf{7.} 6.3\% and 17.5\%.
\textbf{8.} Dark: no capture, $+0.11$ cents of move; lit providers: 0.51 cents of capture, $-0.13$ of move.
\textbf{9.} No evidence that block trades in the dark impede \omterm{def:mx:fragmentation-and-routing:discovery}{price discovery}.
\textbf{10.} The pinger probes every two seconds with 100 shares and trades 500 ahead after two fills in a row; 30\% of block sellers never firm up and buy
2\,000 ahead.
\textbf{11.} See exercise~6: 4.95 basis points.
\textbf{12.} $\lambda\times$ each triggered trade ahead $\times$ the fund's shares still to fill, over the order's size.
\textbf{13.} Lit 3.8, 7.3, 4.95 own impact; dark 12.5, 26.1, 10.9 leakage; dark with a minimum 9.5, 11.9, 1.4 own impact; block $-2.9$, $-1.7$, 0.6
leakage (standard errors 0.5 to 2.3).
\textbf{14.} The pinger found it 88 times and bought 44\,000 shares ahead: the leak (10.9 basis points) is larger than the lit order's own impact (4.95).
\textbf{15.} It stops the leak; only 46\% fills in the dark and the rest is bought at the end.
\textbf{16.} Their sellers bring size that would not otherwise have traded; with leaky or informed counterparties, or when the block's sellers would
otherwise have sold on the lit book, they would not be.
\textbf{17.} ITG: an undisclosed desk traded against its dark pool's subscribers using their live orders (USD 20.3 million). Barclays: misrepresented how
it policed its pool (USD 70 million). Credit Suisse: misrepresented its scoring, took sub-penny orders and alerted high-frequency firms to customer
orders (USD 84.3 million).
\textbf{18.} The double volume cap was replaced by a single volume cap on the reference price waiver, from 18 months after the MiFIR review's entry into
force.
\textbf{19.} \emph{Named result}: over the 30-minute window the mid drifted 26.1 basis points against the unprotected dark order and 7.3 against the lit
one, and the dark order's accounted leakage cost 10.9 basis points; a minimum quantity above the probe size removed it.
\textbf{20.} To block venues with natural sellers first, and to the dark pool only behind a minimum quantity larger than any probe.
\end{solution}

\begin{solution}{iq:mx:dark-pools-and-blocks:1}
A venue that displays no orders and usually crosses at the lit midpoint: an institution uses it to trade size without showing it, saving the spread and
hoping to avoid impact.
\iqlookfor{Midpoint; no display; size and leakage.}
\end{solution}

\begin{solution}{iq:mx:dark-pools-and-blocks:2}
A \omterm{def:mx:dark-pools-and-blocks:minqty}{minimum execution quantity} above typical probe sizes, randomised and partial exposure across venues, avoiding venues with poor participant controls,
conditionals only on venues that police \omterm{def:mx:dark-pools-and-blocks:firmup}{firm-up rates}, and monitoring fills for signs of probing.
\iqlookfor{Minimum quantity; venue selection; monitoring.}
\end{solution}

\begin{solution}{iq:mx:dark-pools-and-blocks:3}
Theory (Zhu): dark pools attract uninformed flow and can concentrate information on the exchange. Evidence (Comerton-Forde and Putnins): low levels of
non-block dark trading are benign, high levels harmful, blocks not harmful.
\iqlookfor{Selection argument; nonlinearity in the evidence.}
\end{solution}

\begin{solution}{iq:mx:dark-pools-and-blocks:4}
Compare the price drift while orders are worked with a benchmark of similar orders, the mark-outs of counterparties by venue, and the price moves after
IOIs and invitations; randomise routing where possible to separate leakage from selection.
\iqlookfor{Drift; counterparty mark-outs; a comparison group.}
\end{solution}

\begin{solution}{iq:mx:dark-pools-and-blocks:5}
Per order: owner, side, size, minimum, time, whether it is in an open invitation. Per invitation: both orders, size, start time, who has firmed up. Per
participant: invitations and firm-ups (the rate), and any restriction. Matching in time priority among free orders whose sizes meet both minimums.
\iqlookfor{Invitation state; timer; participant scoring.}
\end{solution}

\begin{solution}{iq:mx:dark-pools-and-blocks:6}
The midpoint saves the spread but the order leaked or took flow out of the lit market, and the price drifted while it was worked; measure the drift and
the counterparties' mark-outs, add a minimum quantity, and compare with lit and block executions of similar orders.
\iqlookfor{Leakage and impact beyond the spread; how to measure them.}
\end{solution}
